% !TEX root = ../main.tex
\chapter{Corrections Log}
\label{app:corrections}

Nine claims were made during this project and later withdrawn. They are listed here in full, each
with the reasoning error named, because the error is more instructive than the number. The canonical
retraction list lives in \texttt{notebooks/cnn\_development\_log.ipynb}; this appendix summarises it
and adds C9, which arose during report writing.

Two failure modes recur, and it is worth naming them before the list. \textbf{C1 and C5 are the same
mistake in opposite directions}: C1 read a trend into noise, C5 read flatness into a range that was
never sampled. Both had recorded the very caveat that should have blocked the conclusion. A stated
caveat is not a substitute for the measurement that resolves it. \textbf{C3, C6 and C9 are premature
conclusions from a single configuration or split}, each welcomed at the time because it made the
result look more credible rather than less.

\section{C1 --- ``\texttt{clutter\_edge} is still climbing past the loss plateau''}

\textbf{Claimed.} A monotone rise across epochs 30, 43 and 50 was the shape of continued
improvement rather than scatter.

\textbf{Corrected.} The three points came from one training trajectory. Seed-to-seed spread at a
\emph{fixed} epoch is $0.054$, nine times the $0.006$ that training from 15 to 30 actually produces.
The reading was as well explained by which seed it was as by how long it trained.

\textbf{Effect on final results.} None. The claim was never carried into a reported number.

\section{C2 --- ``eight seeds will give a standard error of about $0.011$''}

\textbf{Claimed.} When proposing the eight-seed run, extrapolating from a three-seed spread.

\textbf{Corrected.} The achieved standard error was $0.0040$, nearly three times tighter. The
estimate had assumed seed variance is constant across training; it shrinks as a run converges, and
a three-seed spread is a poor variance estimator in any case.

\textbf{Effect on final results.} None --- the error was conservative, so the experiment was better
powered than promised. Recorded because the \emph{mechanism} was misunderstood.

\section{C3 --- ``CFAR genuinely beats the network on \texttt{clutter}''}

\textbf{Claimed.} At \texttt{pos\_weight}~974, the network lost \texttt{clutter} with all eight
seeds below the classical arm, $p = 0.00014$. It was welcomed as the network's genuine weakness.

\textbf{Corrected.} At \texttt{pos\_weight}~100 the same comparison shows no detectable difference.
The deficit was a property of the hyperparameter, not of the detectors. The value 974 had been set
from the measured inverse class frequency and called a principled default, while the log itself
noted that exact class balancing is a choice rather than a derivation.

\textbf{Effect on final results.} The \texttt{clutter} loss is withdrawn. Chapter~\ref{chp:synthetic}
reports that condition as a tie.

\section{C4 --- ``the network beats CFAR on \texttt{isolated}'' ($p = 6.9\times10^{-7}$)}

\textbf{Claimed.} A significant win on the clean homogeneous-noise condition.

\textbf{Corrected.} The $p$-value came from a one-sample test that treats the classical arm's $\Pd$
as exact. It is one draw over 200 frames. With a paired frame bootstrap giving both arms their
sampling error, and Bonferroni across seven conditions, the same comparison gave $p = 0.031$ ---
short of the corrected threshold. Eight seeds of precision on one arm does not license a point
estimate on the other.

\textbf{Effect on final results.} Decisive, and it is why every comparison in this report uses a
paired bootstrap over frames. The final benchmark of Chapter~\ref{chp:synthetic} does establish
\texttt{isolated} under the corrected procedure, on a larger split.

\section{C5 --- ``$\Pd$ is flat with training, so stop at epoch 15''}

\textbf{Claimed.} That $\Pd$ does not depend on training length, so the epoch count could be chosen
on localisation alone.

\textbf{Corrected.} Swept at every epoch, pooled $\Pd$ rises about 26 standard errors from epoch~1
to the plateau. Nothing was measured wrongly: the four original points reproduce exactly. The error
was concluding flatness from four samples that all lay inside the flat region, when the entire rise
sits below the lowest of them.

\textbf{Effect on final results.} The frozen epoch moved from 15 to \TinyEpochs{} under a
pre-registered rule. The difference between them is inside the noise on all three metrics, so the
number was nearly right and the justification was not.

\section{C6 --- ``\texttt{clutter} is a significant loss, so the network does lose somewhere''}

\textbf{Claimed.} On the validation split at the wider gate, a loss clearing the corrected
threshold. Written up as strengthening the position, since a credible learned detector normally
loses somewhere.

\textbf{Corrected.} It did not replicate on held-out data --- a tie, measured with more power, not
less. A result that had just moved across the significance threshold on a single split was treated
as established when the sensible reading was ``unresolved''.

\textbf{Effect on final results.} The final benchmark reports \texttt{clutter} as a tie at the
primary gate, with the difference moving against the network at the wider one but not reaching
significance. The network's real measured weakness is localisation, not detection.

\section{C7 --- ``CFAR's $\alpha$ is calibrated empirically, and the empirical value is used''}

\textbf{Claimed.} As part of the argument that the classical baseline is tuned rather than
inherited.

\textbf{Corrected.} The calibration module was never imported, so all four variants ran on
closed-form multipliers achieving $0.13\times$ to $6.55\times$ the target rate. One variant had
never been in the calibrator's method list at all. The module had been written, validated and
documented, and its \emph{existence} was then treated as evidence that it was in use. \textbf{A
module that exists is not a module that runs.}

\textbf{Effect on final results.} None. Every result is scored at matched \emph{measured} $\Pfa$,
which sweeps the threshold and discards $\alpha$; the defect affected the design operating point
only. Section~\ref{sec:alphacal} now reports the design-versus-measured gap as a finding.

\section{C8 --- the same claim, made true}

\textbf{Resolved.} Multipliers are now Monte-Carlo calibrated on pure noise, verified on a held-out
seed, cached, and read by the detector configuration. The achieved rate moved from
$0.13$--$6.55\times$ target to $0.95$--$1.01\times$.

\textbf{Effect on final results.} None, and this time provably: $\alpha$ is not an input to the
ratio map, so re-deriving every summary with the calibrated table left the detection counts
bit-identical and the cached peaks byte-identical. Recorded as a correction rather than a silent fix
because C7's claim stood in the notebook for weeks before the work behind it was done.

\section{C9 --- the stated justification for the suppression radius}

\textbf{Claimed.} That a suppression radius of \ProtocolNmsOriginal{} had been selected because it
``performed substantially better on isolated targets'' than smaller radii, and that a smaller radius
costs $\Pd$ across every condition by admitting more candidates.

\textbf{Corrected.} The claim was never measured in the form in which it was stated. A controlled
sweep over radii \NmsRadiiTested{}, run on the tuning split and confirmed on the validation split,
finds isolated-target $\Pd$ flat to within $0.005$ across the entire range, for the network and for
every CFAR variant. The mechanism the claim rested on --- candidate density rising sharply as the
radius falls --- is real and large, and is reported in Figure~\ref{fig:nmssel}; the detection
consequence attributed to it is not there once the operating point is rematched. The original
observation appears to have come from a development sweep that watched a single arm.

\textbf{Effect on final results.} Substantial, and unlike C1--C8 it changed reported numbers rather
than only the text. Section~\ref{sec:nmschoice} now reports the full sweep and the criterion the
measurements support. The protocol consequence --- a revised radius, a fresh synthetic holdout and a
re-scored real-data evaluation --- is recorded in Appendix~\ref{app:prereg}. No model was retrained
or fine-tuned, no checkpoint altered, and no threshold derived from anything but synthetic
development data.
